\documentclass[11pt]{article}

\usepackage[margin=1in]{geometry}
\usepackage{amsmath,amssymb,amsthm}
\usepackage{enumitem}

\newtheorem{theorem}{Theorem}
\newtheorem{proposition}{Proposition}

\title{A Lehmann--Romano Proof of Admissibility for Closed Convex Acceptance Regions}
\date{}

\begin{document}
\maketitle

\section{Why exact size at a finite null point is needed}

The statement ``closed and convex acceptance region implies admissibility'' is false if one only assumes
\[
\sup_{\mu\le 0} E_\mu[\phi(X)]\le \alpha.
\]
Indeed, the zero test
\[
\phi(x)\equiv 0
\]
has acceptance region
\[
A_\phi=\mathbf R^k,
\]
which is closed and convex, and it satisfies the size constraint for every \(0<\alpha<1\). But it is not admissible: for example, the test
\[
\psi(x)=1\{x_1>\sqrt{\Sigma_{11}}\Phi^{-1}(1-\alpha)\}
\]
has level \(\alpha\) for testing \(H_0:\mu\le 0\), and has strictly positive power at alternatives with \(\mu_1>0\). Hence \(\psi\) dominates the zero test.

Thus some condition excluding tests of size strictly below \(\alpha\) is necessary. The condition used below is the one in Lehmann and Romano's Corollary 6.7.1: the test has size \(\alpha\) at a finite point \(\mu_0\in\Lambda_0\). This condition is also used directly in the final step of the proof. Once the convexity argument implies that any dominating test must reject wherever \(\phi_0\) rejects, any strict enlargement of the rejection region would force rejection probability strictly above \(\alpha\) at \(\mu_0\), contradicting the level constraint.

\section{A provable admissibility theorem}

\begin{theorem}[Closed convex exact-size tests are \(\alpha\)-admissible]
Let \(0<\alpha<1\). Let
\[
X\sim N(\mu,\Sigma),
\]
where \(\Sigma\) is known and positive definite and \(\mu\in\mathbf R^k\). Consider testing
\[
H_0:\mu\in\Lambda_0=\{\mu:\mu\le 0\}
\quad\text{versus}\quad
H_1:\mu\in\Lambda_1=\mathbf R^k\setminus\Lambda_0.
\]
Let \(\phi_0\) be a nonrandomized test with acceptance region
\[
A_0=\{x:\phi_0(x)=0\}.
\]
Assume:
\begin{enumerate}[label=(\roman*)]
\item \(A_0\) is closed and convex;
\item \(\phi_0\) has level \(\alpha\):
\[
\sup_{\mu\in\Lambda_0} E_\mu[\phi_0(X)]\le \alpha;
\]
\item there exists a finite point \(\mu_0\in\Lambda_0\) such that
\[
E_{\mu_0}[\phi_0(X)]=\alpha.
\]
\end{enumerate}
Then \(\phi_0\) is \(\alpha\)-admissible for testing \(H_0\) against \(H_1\). That is, if \(\phi\) is any level-\(\alpha\) test satisfying
\[
E_\mu[\phi(X)]\ge E_\mu[\phi_0(X)]
\qquad\text{for all }\mu\in\Lambda_1,
\]
then
\[
E_\mu[\phi(X)]=E_\mu[\phi_0(X)]
\qquad\text{for all }\mu\in\Lambda_1.
\]
\end{theorem}

\begin{proof}
We apply Theorem 6.7.1 and Corollary 6.7.1 of Lehmann and Romano, Section 6.7. The relevant result is the following.

\begin{proposition}[Lehmann--Romano convex acceptance region result]
Consider an exponential family with density
\[
p_\theta(t)=C(\theta)\exp\{\theta't\}
\]
with respect to a \(\sigma\)-finite measure \(\nu\), with natural parameter space \(\Omega\). Let \(\Omega_H\) and \(\Omega'\) be disjoint nonempty subsets of \(\Omega\). Let \(\phi_0\) be a nonrandomized test of \(H:\theta\in\Omega_H\), based on \(T\), with acceptance region \(A_0\subseteq\mathbf R^s\). Suppose that \(A_0\) is closed and convex and that the following ray condition holds:

if
\[
A_0\cap\{t:a't>c\}=\varnothing,
\]
then there exist \(\theta^*\in\Omega\) and a sequence \(\lambda_n\to\infty\) such that
\[
\theta^*+\lambda_n a\in\Omega'.
\]
Then any competing acceptance region \(A\) satisfying
\[
P_\theta(T\in A)\le P_\theta(T\in A_0)
\qquad\text{for all }\theta\in\Omega'
\]
is contained in \(A_0\), except for a \(\nu\)-null set. Consequently, the test with acceptance region \(A_0\) is \(d\)-admissible. If the test has size \(\alpha\) and there exists a finite point \(\theta_0\in\overline{\Omega_H}\) such that
\[
E_{\theta_0}[\phi_0(X)]=\alpha,
\]
then the test is also \(\alpha\)-admissible. The same conclusion allows randomized competing tests.
\end{proposition}

We now verify the assumptions of this result in the normal location model.

Write the normal density in canonical exponential-family form. Put
\[
\theta=\Sigma^{-1}\mu,
\qquad
\mu=\Sigma\theta.
\]
With respect to the measure
\[
d\nu(x)=(2\pi)^{-k/2}|\Sigma|^{-1/2}
\exp\left\{-\frac12 x'\Sigma^{-1}x\right\}dx,
\]
the density of \(X\sim N(\Sigma\theta,\Sigma)\) is
\[
p_\theta(x)=\exp\left\{\theta'x-\frac12\theta'\Sigma\theta\right\}.
\]
Thus the sufficient statistic is
\[
T(x)=x,
\]
the natural parameter space is
\[
\Omega=\mathbf R^k,
\]
and the null and alternative in natural-parameter form are
\[
\Omega_H=\{\theta:\Sigma\theta\le 0\},
\qquad
\Omega_K=\{\theta:\Sigma\theta\not\le 0\}.
\]
We apply the Lehmann--Romano result with \(\Omega'=\Omega_K\).

It remains only to verify the ray condition. Let \(a\ne 0\) and \(c\in\mathbf R\) be such that
\[
A_0\cap\{x:a'x>c\}=\varnothing.
\]
Equivalently,
\[
A_0\subseteq\{x:a'x\le c\}.
\]
We claim that
\[
\Sigma a\not\le 0.
\]
Suppose instead that \(\Sigma a\le 0\). For \(\lambda>0\), define
\[
\theta_\lambda=\lambda a,
\qquad
\mu_\lambda=\Sigma\theta_\lambda=\lambda\Sigma a.
\]
Then \(\mu_\lambda\le 0\), so \(\theta_\lambda\in\Omega_H\). Under \(P_{\theta_\lambda}\),
\[
a'X\sim N\left(\lambda a'\Sigma a,\ a'\Sigma a\right).
\]
Since \(\Sigma\) is positive definite and \(a\ne 0\),
\[
a'\Sigma a>0.
\]
Therefore
\[
P_{\theta_\lambda}(a'X>c)\to 1
\qquad\text{as }\lambda\to\infty.
\]
But \(A_0\subseteq\{x:a'x\le c\}\), so
\[
\{x:a'x>c\}\subseteq A_0^c.
\]
Because \(\phi_0\) is nonrandomized, \(\phi_0(x)=1\) on \(A_0^c\). Hence
\[
E_{\theta_\lambda}[\phi_0(X)]
\ge
P_{\theta_\lambda}(a'X>c)
\to 1.
\]
This contradicts the level condition, since every \(\theta_\lambda\in\Omega_H\) and \(\alpha<1\). Thus \(\Sigma a\not\le 0\).

It follows that \(\Sigma a\) has at least one strictly positive component. Therefore, for every \(\lambda>0\),
\[
\Sigma(\lambda a)=\lambda\Sigma a\not\le 0.
\]
Hence
\[
\lambda a\in\Omega_K
\qquad\text{for every }\lambda>0.
\]
The ray condition holds with \(\theta^*=0\) and any sequence \(\lambda_n\to\infty\).

By the Lehmann--Romano result, if \(\phi\) is any level-\(\alpha\) test whose power is at least as large as that of \(\phi_0\) at every alternative, then \(\phi\) must reject wherever \(\phi_0\) rejects, except possibly on a \(\nu\)-null set. Equivalently,
\[
\phi(x)\ge \phi_0(x)
\qquad\text{for }\nu\text{-almost every }x.
\]
Since \(\nu\) has a strictly positive density with respect to Lebesgue measure, this exceptional set has probability zero under every \(N(\mu,\Sigma)\) distribution.

If \(\phi>\phi_0\) on a set of positive Lebesgue measure, then, because the normal density is strictly positive everywhere,
\[
E_{\mu_0}[\phi(X)]>E_{\mu_0}[\phi_0(X)]=\alpha.
\]
This contradicts the level-\(\alpha\) constraint, because \(\mu_0\in\Lambda_0\). Therefore
\[
\phi=\phi_0
\qquad\text{almost everywhere}.
\]
Consequently,
\[
E_\mu[\phi(X)]=E_\mu[\phi_0(X)]
\qquad\text{for all }\mu\in\Lambda_1.
\]
This proves \(\alpha\)-admissibility.
\end{proof}

\begin{thebibliography}{9}

\bibitem{LehmannRomano2005}
Lehmann, E. L. and Romano, J. P. (2005).
\textit{Testing Statistical Hypotheses}.
Third edition. Springer.

\end{thebibliography}

\end{document}
